
\begin{corollaire}
\label{XIII.5.7}
Las notaciones son las de~\Ref{XIII.5.6}. Sean $\overline{S}$ la
localización estricta de~$S$ en~$\overline{x}$,
$$
\overline{g}\colon\overline{U}\to\overline{S}
\quad \text{y} \quad
\tilde{g}\colon\widetilde{U}\to
\ifthenelse{\boolean{orig}}
{\widetilde{S}}
{\overline{S}}
$$
los morfismos canónicos. Entonces los morfismos $\overline{g}$ y
$\tilde g$ son $0$-acíclicos \textup{(SGA~4 XV~1.3)}. Sean $G$ un haz
de grupos constructible sobre~$\overline{S}$,
$F=\overline{g}^*G$ y $P$ un torsor bajo~$F$. Para que $P$ esté
moderadamente ramificado sobre~$\overline{X}$ relativamente
a~$\overline{S}$, es necesario y suficiente que su imagen inversa
$\widetilde{P}$ sobre~$\widetilde{U}$ sea trivial.
\end{corollaire}

En efecto, para todo esquema
$\overline{X'}=\overline{X}[T_i]_{i\in I(x)}\big/
(T_i^{n_i}-f_i)$, donde los $n_i$ son enteros $>0$ primos con~$p$, el
morfismo $\overline{f}'\colon\overline{X'}\to\overline{S}$ es
$0$-acíclico. Por tanto, las fibras geométricas de~$\overline{f}'$ en
los distintos puntos de~$\overline{S}$ son conexas e incluso
irreducibles. Lo mismo sucede con las fibras geométricas de los
morfismos $\overline{g}'\colon\overline{U'}\to\overline{S}$, lo que
demuestra que los $\overline{g}'$, y por tanto también~$\tilde g$, son
$0$-acíclicos (SGA~4 XV~1.16).

Es claro que un torsor $P$ sobre~$\overline{U}$, bajo~$F$, cuya imagen
inversa sobre~$\widetilde{U}$ es trivial, está moderadamente ramificado
sobre~$\overline{X}$ relativamente a~$\overline{S}$. Demostremos el
recíproco: si $P$ está moderadamente ramificado sobre~$\overline{X}$
relativamente a~$\overline{S}$, su imagen inversa sobre~$\widetilde{U}$
es trivial.

Se deduce de (SGA~4 IX~2.14~(ii)) que pueden encontrarse un morfismo
finito $n\colon S_1\to\overline{S}$, un haz de grupos constante~$C$
sobre~$S_1$ y un monomorfismo $G\to n_*C$. Consideremos el siguiente
diagrama conmutativo formado por cuadrados cartesianos:
$$
\xymatrix{ \widetilde{U}_{1} \ar[r] \ar[d]_{r} &U_{1} \ar[r]^{g_{1}}
\ar[d]_{q} &S_{1} \ar[d]_{n} \cr \widetilde{U} \ar[r] &\overline{U}
\ar[r]^{\overline{g}} &\overline{S} \rlap{\quoi.}\cr }
$$
Sean
\marginpar{435}
$C_1$ (\resp $\widetilde{C}_1$) la imagen inversa de~$C$ sobre~$U_1$
(\resp sobre~$\widetilde{U}_1$). Se tiene un diagrama conmutativo en el
que $i$ y $j$ son isomorfismos (SGA~4 VIII~5.8):
\begin{equation*}
\label{eq:XIII.5.7.*}
\tag{$*$}
\begin{array}{c}
\xymatrix{ \H^{1}(\overline{U},q_*C_{1}) \ar[r]^{i}
\ar[d] &\H^{1}(U_{1},C_{1}) \ar[d] \cr
\H^{1}(\widetilde{U},r_*\widetilde{C}_{1}) \ar[r]^{j}
&\H^{1}(\widetilde{U}_{1},\widetilde{C}_{1}) \rlap{\quoi.}\cr }
\end{array}
\end{equation*}
Sea $Q$ el torsor bajo~$q_*C_1$ deducido de~$P$ por extensión del grupo
estructural $F\to q_*C_1$. Según~\Ref{XIII.2.1.4}, $Q$ está
moderadamente ramificado sobre~$\overline{X}$ relativamente
a~$\overline{S}$. Al torsor~$Q$ corresponde, mediante~$i$, un torsor
$Q_1$ bajo~$C_1$, y es claro que $Q_1$ está moderadamente ramificado
sobre $X_1=\overline{X}\times_{\overline{S}}S_1$ relativamente a~$S_1$.
Se deduce entonces de~\Ref{XIII.5.6} que la imagen inversa
$\widetilde{Q}_1$ de~$Q_1$ sobre~$\widetilde{U}_1$ es trivial, y el
diagrama~\eqref{eq:XIII.5.7.*} demuestra que la imagen inversa
$\widetilde{Q}$ de~$Q$ sobre~$\widetilde{U}$ es trivial.

Consideremos el siguiente diagrama conmutativo, cuya segunda fila es
exacta (SGA~4 XII~3.1):
\begin{equation*}
\label{eq:XIII.5.7.**}
\tag{$**$}
\begin{array}{c}
\xymatrix{ \H^{0}(\overline{S},n_*C/G) \ar[r]
\ar[d]_{k} &\H^{1}(\overline{S},G)\rlap{$~=~1$} \ar[d] \cr
\H^{0}(\widetilde{U},r_*\widetilde{C}_{1}/\widetilde{F}) \ar[r]
&\H^{1}(\widetilde{U},\widetilde{F}) \ar[r]
&\H^{1}(\widetilde{U},r_*\widetilde{C}_{1}) \rlap{\quoi.} \cr }
\end{array}
\end{equation*}
Como el morfismo $\widetilde{U}\to\overline{S}$ es $0$-acíclico,
$k$ es un isomorfismo. De~\eqref{eq:XIII.5.7.**} se deduce entonces
que $\widetilde{P}$ es trivial.
